The Methods follow the same evidentiary sequence as the Results. We first define cohorts,
labels, analysis units and frozen representations; then specify cross-cohort and
patient-disjoint evaluation; and finally describe confounder, site, setting and endpoint
audits, uncertainty and provenance. This ordering keeps each qualification decision tied to
the data and analysis that can support it.

\subsection{Cohorts and labels}

This subsection defines the data source, label provenance, analysis unit, and missing-outcome
rule for every target before any model evaluation is specified.

NADT-Prostate provided paired prostate histology, per-core or per-slide grade and phenotype
labels, and patient identifiers for source-probe fitting and patient aggregation. PANDA
provided institution-labelled case images and official ISUP grades; its tumor/benign
phenotype was derived as ISUP 0 versus ISUP at least 1. TCGA-PRAD provided histology,
patient identifiers, grade, PTEN copy-number loss, SPOP mutation and AR-activity labels, as
well as separately constructed recurrence-style endpoints. PRECISE provided expert-annotated
serial sections and per-image pathologist Gleason scores. LEOPARD provided biochemical
recurrence event and follow-up labels for recurrence-model fitting. The public NADT,
PANDA and PRECISE resources are described in their source records
\cite{nadt2021,panda2022,precise2026}; LEOPARD is described in the Medical Imaging with Deep
Learning (MIDL) 2026 challenge
report \cite{leopard2026}, and the organizer baseline remains an arXiv preprint
\cite{grisi2026}. Each analysis used only
rows with an evaluable target; missing or uncertain outcomes were not converted to negative
labels.

For NADT patient-level evaluation, labels and predictions were aggregated within patient;
the phenotype target was the patient's mean tumor fraction rather than a rounded binary
label. PANDA sampling capped each institution--ISUP stratum at 100 cases before tissue
filtering; its grade-derived phenotype was not a separately supplied annotation.

The recurrence endpoints were retained under their source definitions. The Reconstructed
with-tumor endpoint used a recorded with-tumor follow-up as an event and censored at the
latest valid non-event follow-up. The strict recurrence-only sensitivity required a
tumor-free visit before a later with-tumor visit. TCGA-CDR progression-free survival (PFS)
and disease-free survival (DFS) used their named cBioPortal comparator fields. Official
TCGA-CDR PFI used the official PFI event and time
fields \cite{tcgacdr2018}. No
endpoint substitution was renamed as official PFI.

\subsection{Frozen embeddings and aggregation}

With the cohorts defined, we describe how frozen image representations were constructed and
aligned to the analysis unit required by each endpoint.

CONCH and Virchow were used as frozen encoders; encoder weights were not fine-tuned.
Tissue-qualified tiles were embedded and averaged to slide representations. Where the
endpoint unit was the patient, multiple slide values were aggregated by patient before
primary scoring. PANDA used case-image rows, PRECISE used imaging sessions, AR site effects
used slide-level estimates with patient-cluster resampling, and all such departures from a
patient-level effect unit were recorded in the figure source tables.

The sensitivity analysis crossed each encoder with native and shared physical scales, tile
budgets of 16, 32 and 64, and five repeated sampling seeds. These settings were applied to
the fixed cohorts and folds. They were not treated as newly sampled populations.

\subsection{Fixed probes and patient-disjoint evaluation}

This subsection specifies how the fixed representations were mapped to targets while
preventing information from the same patient from entering both training and evaluation.

Continuous marker-family targets used standardized ridge regression with the saved fixed
alpha grid. Binary TCGA-PRAD screening tests used standardized logistic regression with
fixed regularization and balanced class weighting. The frozen NADT phenotype probe used for
PANDA transfer was instead \texttt{LogisticRegression(C=1.0)} with 2,000 maximum iterations
and without class weighting, matching its saved generating script. Within-cohort evaluation
used five patient-disjoint grouped folds, keeping all slides from a patient within one fold.
Continuous targets were scored by patient Spearman correlation, binary targets by patient
AUROC, and time-to-event targets by patient-level concordance and the saved time-dependent
metrics.

For zero-shot transfer, the source-cohort probe and preprocessing were applied to the target
cohort without refitting. Native refits and held-out site analyses were reported as distinct
evaluation types and did not substitute for zero-shot transfer.

\subsection{Qualification criteria and multiplicity}

We next define how individual estimates were converted into evidence states and how the
saved hypothesis family was controlled for multiplicity.

Transportable evidence required an appropriate cross-cohort evaluation with consistent
direction. Context-sensitive evidence indicated material dependence on site, encoder,
physical scale, tile budget or endpoint. Unsupported evidence indicated that the frozen
primary analysis did not support the proposed effect; it was not interpreted as evidence of
biological absence. These dimensions address published concerns about molecular-label
confounding, task-dependent foundation-model rankings and institution-associated technical
signals \cite{dawood2026,neidlinger2025,pathorob2026}.

The saved 17-test family comprised 13 marker hypotheses---the six initial candidates and
seven additional screening candidates---plus four nested/refit confounder audits: PTEN
beyond grade, AR beyond grade, and the post-hoc recurrence risk signal beyond grade or a
fully adjusted clinical model.
Continuous marker hypotheses used two-sided Spearman tests; binary marker hypotheses used
two-sided Mann--Whitney tests. The fully refitted, grade-stratified audit permutations used
one-sided refit-permutation $p$ values for positive improvement. Statistical support used
$\alpha=0.05$ with Benjamini--Hochberg false-discovery-rate adjustment across this family.
Re-tests of an existing hypothesis in another cohort or encoder were qualification evidence
and were not added to the discovery family.

\subsection{Confounder and site analyses}

These analyses test two alternative explanations for pooled performance: information already
captured by grade and effects that do not remain precise within source sites.

PTEN and AR increment analyses compared clinical-only, image-only and combined models in
outer patient-disjoint folds. Inner out-of-fold image predictions were used to construct
outer-training meta-model features, and the image probe was refitted on each complete
outer-training fold before application to untouched outer-test patients. Primary increments
were the paired held-out change in AUROC for PTEN and change in $R^2$ for AR.

The nonparametric check permuted patient labels within rounded patient-mean Gleason strata
and refitted the full nested pipeline. AR site analyses estimated slide-level Spearman
correlations and resampled patients as clusters. Site descriptors and pooled estimates were
kept at their recorded units. Raw scanner identifiers were not interpreted as validated
scanner models, and absent stain metadata were not inferred.

\subsection{Correlated stability analysis}

This subsection evaluates dependence on encoder, physical scale, tile budget, and sampling
seed without treating reused cohorts and folds as independent validation data.

For each marker, seed-level scores were summarized into encoder--scale--tile configurations.
The saved summaries report mean, sample standard deviation, range, chance-or-worse count and
a Student-$t$ interval across the five tile-sampling seeds. AUROC and C-index used a null of
$0.5$; Spearman correlation used zero. Paired comparisons preserved the seed-level pairing
for scale, encoder and tile-budget contrasts.

Because cohorts, patients and folds were shared across settings, all grid summaries were
interpreted as correlated sensitivity analyses. Seed intervals were descriptive sampling
intervals, not patient-bootstrap or population-confidence intervals. No grid setting was
used to revise frozen labels, scores, folds or evidence thresholds.

\subsection{Recurrence endpoints and paired comparisons}

We then define the exploratory recurrence signal, keep the competing endpoint definitions
separate, and specify the paired comparisons used to assess incremental value.

The post-hoc recurrence risk signal was a Cox risk score fitted in LEOPARD and applied
without target-cohort refitting to TCGA-PRAD. Its existing principal-component analysis
(PCA) dimension sweep served only as exploratory sensitivity analysis. PCA-dimension
selection was not nested within an outer evaluation,
so the sweep did not provide an unbiased model-selection estimate and no dimension was
interpreted as an unbiased optimum.

Paired TCGA-PRAD comparisons used the same complete-case patients and the same bootstrap
draw for the clinical model and its image-augmented counterpart. Models were compared for
the Reconstructed with-tumor endpoint and Official TCGA-CDR PFI separately. C-index
improvement was defined as augmented minus reference C-index; integrated-Brier-score
improvement was reference minus augmented IBS, so positive values favored the augmented
model for both metrics. Reconstructed recurrence, PFS, DFS and official PFI were never
pooled into one endpoint.

\subsection{Bootstrap uncertainty and missingness}

This subsection states how uncertainty was estimated and how unavailable outcomes or
unretained bootstrap accounting were handled without reconstruction or recoding.

Where the saved source supported this construction, patient-level estimates used
95\% percentile patient-bootstrap intervals. AR site rows used a patient-cluster bootstrap
around a slide-level effect. For legacy NADT/TCGA primary and AR site interval summaries,
saved tables retained interval endpoints but not requested, valid and undefined accounting,
except requested 2,000 where explicitly recorded; unavailable counts were not reconstructed.
Paired recurrence-risk outputs separately retained 2,000 valid and zero undefined replicates.
Paired survival contrasts used identical patient draws for both models. Outputs with
replicate accounting retained and reported undefined replicate counts and fractions;
estimates without a saved interval were marked unavailable. Analyses used their
saved evaluable or complete-case frames; no missing, uncertain or unevaluable value was
recoded as a negative outcome. Multiple imputation was not performed for the recurrence-risk paired
comparison.

\subsection{Software and provenance}

Finally, we describe the records that connect each manuscript result to its code, saved
source, execution configuration, and immutable clinical-review input.

Analyses were executed with the repository's locked Python environment and auditable entry
scripts. Fixed random seeds, software versions, input hashes, endpoint mappings and output
hashes were recorded in saved run configurations and manifests. Main figures were rendered
from saved comma-separated-value (CSV) source tables rather than transient objects.
Source-table schema, analysis
unit, denominator, endpoint identifier and interval type were checked before a result was
used in the manuscript. The clinician review source for PRECISE remained immutable throughout
the manuscript work.

\paragraph{Large language model assistance.}
An assistive large language model, OpenAI Codex, was used during code and manuscript
revision for code editing, test authoring, figure-production assistance and manuscript
language revision. It was not treated as an author and did not independently determine the
study endpoints or scientific conclusions. The exact product/model version, usage dates and
accountable authors' final verification statement require author confirmation before
submission.
